\section{Introduction}

The volume of genomic variant data is growing rapidly. Population-scale
sequencing projects now routinely produce Variant Call Format (VCF) files that
describe thousands of samples and millions of variant sites per chromosome
\citep{vcftools}. Once variants have been called, these files are scanned
repeatedly for quality filtering, cohort subsetting, population-genetic
summaries, and conversion to downstream formats. These repeated scans are a
recurring cost in post-calling analysis, and each is typically carried out by a
tool that parses the entire record even when the operation depends on only one or
two of its fields.

This inefficiency follows directly from the VCF design. A VCF record stores
every property of a variant on a single line, beginning with the site-level
fields (position, quality, filter status, and INFO annotations) and continuing,
for every sample in the cohort, with a FORMAT block such as
\texttt{GT:DP:GQ:AD} that encodes the genotype, read depth, genotype quality,
and allelic depths. This layout makes VCF a general exchange format that any
downstream program can read, which is why it has become the standard
representation for variant data. It also means that a program answering a
field-limited question, such as retaining sites with $\mathtt{QUAL}>30$ or
measuring per-site missingness, must still decode the full record, including the
per-sample values it will not use. On cohorts of thousands of samples, where the
per-sample FORMAT block is by far the largest part of each record, decoding these
unused fields typically dominates the running time of the operation.

Established tools take different routes to this problem. VCFtools
\citep{vcftools} parses the full text record. bcftools with HTSlib
\citep{bcftools,htslib} unpacks BCF fields in tiers and, given a tabix or CSI
coordinate index \citep{tabix}, restricts reading to requested genomic regions,
so it avoids decoding data it does not need for site-level BCF queries. Fast
parsing libraries such as cyvcf2 \citep{cyvcf2} likewise provide lazy access to
individual fields, and dedicated engines re-encode genotypes into
query-optimised stores for fast field- or sample-limited access, including GQT
\citep{gqt} and the fixed-schema binary genotypes of PLINK \citep{plink2}. These
approaches are effective within their scope, but none provides predicate-driven
decoding of unmodified VCF text---where the single-line layout otherwise forces a
full parse---combined with block-level skipping keyed on non-positional
predicates, such as filter status or allele-frequency thresholds, that a
coordinate index cannot exploit.

VariantFlow addresses this inefficiency through selective execution. For each
filter, statistic, or projection, it determines the minimal set of fields on
which the result depends and decodes only those fields while streaming once
through the file, leaving the remainder of each record unparsed; where the
operation allows, it returns the original record text unchanged. VariantFlow
applies this strategy to VCF/BCF filtering, VCFtools-compatible
population-genetic summaries, and columnar export, and it reports a performance
gain only where its output matches an established reference tool. It is intended
to complement, not replace, bcftools, HTSlib, VCFtools, and GATK \citep{gatk}.
